\cotefile{29}
\scanpage{14}
Appelons \uline{crible couvrant} de $\underline{C}$ toute sous-catégorie \inserted{\unsure{\deleted{n}p|\deleted{p}n}\unsure{\deleted{e}l|\deleted{l}e|\deleted{d}l|\deleted{l}d}eine} $\underline{C}_o$ de $\underline{C}$ qui est un crible (i.e. pour toute flèche, l'origine appartient à $\underline{C}_o$ s'il en est de même de l'extrémité) et tel que les $\underline{p}(M)$ pour $M\quad\mathrm{Ob}\,\underline{C}_o$ recouvrent l'objet final $S$ de $\underline{\mathrm{Et}}_S$. Pour tout $\underline{C}_o$, soit $\mathrm{Rev}^{\underline{R}}_{\underline{C}_o}(S)$ la catégorie des sections cartésiennes de $\underline{D}$ au d essus de $\underline{C}_o$. Pour $\underline{C}_o$ var\unsure{\deleted{s}i|\deleted{i}s}able\deleted{\unsure{s}}, ces catégories forment \deleted{\unsure{un|on|ou} les fibres d'une}\inserted{un systême inductif de} catégories\deleted{fibrée\unsure{s} sur la catégorie ordonnée de ces $\underline{C}_o$} (NB si $\underline{C}'_o\quad\underline{C}_o$, on obtient un foncteur restriction évident $\mathrm{Rev}^{\underline{R}}_{\underline{C}_o}(S)\quad\mathrm{Rev}^{\underline{R}}_{\underline{C}'_o}(S)$), \deleted{et} dont la limite inductive est notée $\mathrm{Rev}^{\underline{R}}(S)$ et appelée \uline{catégorie des revêtements de} $S$ \uline{de type} $\underline{R}$.

Signaons que dans tous les cas que nous aurons à utiliser, les foncteurs de transition $\mathrm{Rev}^{\underline{R}}_{\underline{C}_o}(S)\quad\mathrm{Rev}^{\underline{R}}_{\underline{C}'_o}(S)$ sont pleinement fidèles, donc aussi les foncteurs $\mathrm{Rev}^{\underline{R}}_{\underline{C}_o}(S)\quad\mathrm{Rev}^{\underline{R}}(S)$, qui permettent donc d'identifier (à équivalence près) les $\mathrm{Rev}^{\underline{R}}_{\underline{C}_o}(S)$ à des sous-catégories pleines de la catégorie $\mathrm{Rev}^{\underline{R}}(S)$ des revêtements detype $S$. Nous supposerons \deleted{de plus} que pour tout $\underline{C}_o$, \deleted{le foncteur composé} et tout $X\quad\mathrm{Ob}\,\mathrm{Rev}^{\underline{R}}_{\underline{C}_o}(S)$, le foncteur composé $\underline{C}_o\overset{X}{\qquad}\underline{D}\overset{\mathrm{source}}{\qquad}(\mathrm{Sch})_{/S}$, (où le deuxième foncteur est le foncteur déduit du foncteur source sur la catégorie des s chémas finis étales relatifs,) admet une limite inductive dans $(\mathrm{Sch})_{/S}$, que nous noterons $\unsure{\deleted{r}\underline{r}|\underline{r}|\deleted{f}r|\deleted{r}f}(X)$, et que pour $C'_o\quad\underline{C}_o$, désignant par $X'$ la restriction de $X$ à $\underline{C}'_o$, l'homomorphisme canonique $r(X'_o)\quad r(X)$ est un isomorphisme. (Ces conditions seront satisfaites dans les exemples \inserted{(de généralité croissante)} 1 à 3 ci-dessous). Les différents foncteurs $r$ se r accorde\unsure{nt|\deleted{m}nt|\deleted{nt}m} alors et
% archive sheet number: 13
